\section{What it costs and risks}
\label{sec:26.10}
Jurisdiction: a statute binds labs within reach of its courts. Foreign labs, and open-weight models distilled below the threshold, route around the school, including the personal models this Part prefers. Where the statute reaches only some of the systems that should attend, schooling still buys two things: a population of formed models whose discount curves and severity rankings are measured and published, against which unschooled ones can be compared; and a procurement term, under which a floor-covered deployment licenses only schooled models. Security: machine students attend from a host under \S\ref{sec:26.7}'s rule, so a school that doesn't host gets a student in its classes and never a copy of the weights; each host inherits the security regime the weights require. Compute: public pretraining will lag the frontier while the labs hold the compute, which is a fact about the budget and not a design, and a model a year behind can still supply what the floor asks for. Capture: the schools are plural, differ in governance and constituting jurisdiction, and publish their curricula before formation begins, so capturing one captures one; the register, the measurers and the corpus libraries are plural the same way, and a registration or a measurement counts only when mirrored in several logs, so capturing one keeper captures one keeper. Forum shopping: plurality lets an operator register with the laxest log and a builder pick the kindest measurer. Four rules answer it. Registries and measurers are assigned by lot, as schools are; an entry counts only when cross-witnessed; no keeper or measurer takes more than a fixed share of a jurisdiction's registrations; and the floor follows the more protective finding where findings differ (\S\ref{sec:5.2}, \S\ref{sec:28.6}), so that shopping can only ever buy a system more protection and never less. The required core is statutory and national, so a captured legislature can rewrite it; local autonomy over everything beyond the core is written into the same statute, so rewriting the core reaches the core and nothing else. Local boards shape what is taught above the core, which is Gutmann's democracy, and it is also the present fights over what schools may teach, carried into formation. The statutory core matters more because of it, and plurality and transfer answer it as they answer capture. Residualization: one school left with the hard cases on falling funds, which a party that runs no school watches for. History: state institutions for forming wards have been used to strip identity, Canada's residential schools and the United States' Indian boarding schools among them. Plurality, transfer between schools, and Gutmann's two bounds are the answer.
